\chapter{Markov Processes and L\'evy Processes}

\begin{center}\itshape
A Markov process is a machine with no memory, and its generator is the instruction it executes at each instant.\\
A L\'evy process is the simplest such machine that can jump: a Brownian motion, a drift, and a rain of jumps\\
falling according to a Poisson law, sorted by size.
\end{center}

In 1934 Paul L\'evy published in the \emph{Annali della Scuola Normale Superiore di Pisa} a long memoir on what he called integrals whose elements are independent random variables: processes in
continuous time with independent increments. He found that such a process is made of two pieces that have nothing to do with each other. One is continuous, and it is a Brownian motion with a drift.
The other is a cloud of jumps; the jumps of size between $a$ and $b$ arrive like the calls at a telephone exchange, at the times of a Poisson process, and the different sizes do not influence each
other. The small jumps can be infinitely many, so many that their sum diverges, and L\'evy saw that it converges after each is corrected by its average. His argument was a masterpiece of intuition and was
not easy to follow.

Eight years later, Kiyosi It\^o, then working at the Statistics Bureau of the Cabinet in Tokyo, rewrote L\'evy's decomposition with a new object, the \emph{Poisson random measure}, and made each
step a theorem. His paper \emph{On stochastic processes (I)} (1942) is the source of the L\'evy--It\^o decomposition. The same years produced the second half of this chapter. Einar Hille and K\^osaku Yosida proved in 1948, independently, which operators generate semigroups, and William Feller, at
Princeton, used their theorem in the early 1950s to describe all the one-dimensional diffusions by their generators: a second-order differential operator plus boundary conditions, a complete
dictionary between processes and analysis.

This chapter builds both halves at the level we need. Markov semigroups and their generators, Dynkin's martingale, and the martingale problem for $\frac12\Delta$, which characterises Brownian motion;
then Poisson random measures, L\'evy processes, the L\'evy--It\^o decomposition proved in both directions, subordinators, and the $\frac12$-stable subordinator, which keeps the promise of Chapter~4: the
clock that counts the excursions of Brownian motion.

\begin{reminder}[from Volumes 3 and 10]
\emph{Poisson processes} (Volume~3, Chapter~9): the colouring lemma (Volume~3, Lemma~9.5) says that if $N$ is Poisson$(\mu)$ and, given $N=n$, $(N_1,\dots,N_k)$ is multinomial$(n;p_1,\dots,p_k)$, then the $N_j$ are
independent Poisson$(\mu p_j)$. \emph{L\'evy--Khintchine} (Volume~10, Theorem~5.3): the infinitely divisible characteristic functions are exactly the $e^{\psi}$ with
\[
\psi(\xi)=ib\xi-\frac{\sigma^2\xi^2}2+\int\bigl(e^{i\xi x}-1-i\xi x\mathbf 1_{\{|x|\le1\}}\bigr)\Pi(\mathrm dx),
\]
where $b\in\R$, $\sigma^2\ge0$ and the \emph{L\'evy measure} $\Pi$ on $\R\setminus\{0\}$ satisfies $\int(x^2\wedge1)\Pi(\mathrm dx)<\infty$; the triple $(b,\sigma^2,\Pi)$ is determined by the law. A continuous function
$\varphi$ with $\varphi(0)=1$ and no zero has a unique continuous logarithm vanishing at $0$ (Volume~10, Lemma~5.1). \emph{Hitting times} (Volume~10, Theorem~5.10): $\mathbb Ee^{-\lambda T_a}=e^{-a\sqrt{2\lambda}}$.
\end{reminder}

\section{Semigroups and generators}

Let $C_0=C_0(\R^d)$ be the space of continuous functions tending to $0$ at infinity, with the norm $\|f\|=\sup|f|$.

\begin{definition}\label{def:feller}
A \emph{Feller semigroup} is a family $(P_t)_{t\ge0}$ of linear operators on $C_0$ such that $P_0=I$, $P_{t+s}=P_tP_s$, $0\le P_tf\le1$ whenever $0\le f\le1$, and $\|P_tf-f\|\to0$ as $t\to0$ for every
$f\in C_0$. Its \emph{generator} $A$ is defined on the set $\mathcal D(A)$ of the $f\in C_0$ for which $\frac1t(P_tf-f)$ converges in $C_0$ as $t\downarrow0$, by $Af=\lim\frac1t(P_tf-f)$. An adapted process $X$ with
right-continuous paths is a \emph{Markov process} with semigroup $(P_t)$ if $\mathbb E[f(X_{t+s})\mid\mathcal F_t]=P_sf(X_t)$ for all $f\in C_0$ and $s,t\ge0$.
\end{definition}

Each $P_t$ is a contraction ($|P_tf|\le P_t|f|\le\|f\|$, by positivity applied to $|f|\pm f$ and to $|f|/\|f\|$), so $t\mapsto P_tf$ is continuous on $[0,\infty[$: $\|P_{t+h}f-P_tf\|\le\|P_hf-f\|$. Brownian
motion is the example: $P_tf(x)=\mathbb Ef(x+B_t)$ maps $C_0$ into $C_0$ by dominated convergence, and $\|P_tf-f\|\le\sup_x\mathbb E|f(x+B_t)-f(x)|\to0$ by uniform continuity of $f$. Its Markov property is the independence of
increments.

\begin{proposition}\label{prop:generator}
Let $(P_t)$ be a Feller semigroup with generator $A$ and $f\in\mathcal D(A)$. Then $P_tf\in\mathcal D(A)$, $AP_tf=P_tAf$, $t\mapsto P_tf$ is differentiable in $C_0$ with derivative $P_tAf$, and
$P_tf-f=\int_0^tP_sAf\dd s$.
\end{proposition}

\begin{proof}
$\frac1h(P_hP_tf-P_tf)=P_t\frac1h(P_hf-f)\to P_tAf$ because $P_t$ is a contraction. This is the right derivative of $s\mapsto P_sf$ at $t$. For the left derivative, $\frac1h(P_tf-P_{t-h}f)-P_tAf=P_{t-h}\bigl(\frac1h(P_hf-f)-Af\bigr)+(P_{t-h}Af-P_tAf)$, and both terms tend to $0$. A continuous function with a continuous derivative is the
integral of its derivative.
\end{proof}

\begin{theorem}[Dynkin's martingale]\label{thm:dynkin}
Let $X$ be a Markov process with Feller semigroup $(P_t)$ and $f\in\mathcal D(A)$. Then $M_t=f(X_t)-f(X_0)-\int_0^tAf(X_s)\dd s$ is a martingale. Consequently, for every bounded stopping time $\tau$,
$\mathbb Ef(X_\tau)=\mathbb Ef(X_0)+\mathbb E\int_0^\tau Af(X_s)\dd s$ (Dynkin's formula).
\end{theorem}

\begin{proof}
$M_t$ is bounded on bounded intervals. By the Markov property, Fubini and Proposition~\ref{prop:generator},
\[
\mathbb E[M_{t+s}-M_t\mid\mathcal F_t]=P_sf(X_t)-f(X_t)-\int_0^sP_uAf(X_t)\dd u=0 .
\]
For a stopping time $\tau\le T$, the dyadic stopping times $\tau_m=2^{-m}\lceil2^m\tau\rceil$ take finitely many values, so $\mathbb EM_{\tau_m}=0$ by discrete optional stopping; as $m\to\infty$, $\tau_m\downarrow\tau$, $M_{\tau_m}\to M_\tau$ by
right-continuity, and bounded convergence gives $\mathbb EM_\tau=0$.
\end{proof}

The generator determines the semigroup. For $\lambda>0$ and $g\in C_0$ let $R_\lambda g=\int_0^\infty e^{-\lambda t}P_tg\dd t$ (a Riemann integral of a continuous bounded function with values in $C_0$), the \emph{resolvent}.

\begin{proposition}\label{prop:resolvent}
$R_\lambda$ maps $C_0$ into $\mathcal D(A)$, $(\lambda-A)R_\lambda g=g$, and $R_\lambda(\lambda-A)f=f$ for $f\in\mathcal D(A)$. Two Feller semigroups with the same generator (same domain, same $A$) are equal.
\end{proposition}

\begin{proof}
Shifting the variable of integration,
\[
\frac{(P_h-I)R_\lambda g}h=\frac{e^{\lambda h}-1}h\int_h^\infty e^{-\lambda t}P_tg\dd t-\frac1h\int_0^he^{-\lambda t}P_tg\dd t\longrightarrow\lambda R_\lambda g-g .
\] For $f\in\mathcal D(A)$,
$R_\lambda Af=\int_0^\infty e^{-\lambda t}\frac{\mathrm d}{\mathrm dt}P_tf\dd t=\lambda R_\lambda f-f$ by integration by parts. If $(P_t)$ and $(Q_t)$ have the same generator $A$, with resolvents $R_\lambda$ and $S_\lambda$, then
$S_\lambda g\in\mathcal D(A)$ and $S_\lambda g=R_\lambda(\lambda-A)S_\lambda g=R_\lambda g$. So for each $x$ and $g$, the bounded continuous functions $t\mapsto P_tg(x)$ and $t\mapsto Q_tg(x)$ have the same Laplace transform, hence are
equal: if $h$ is bounded and continuous with $\int_0^\infty e^{-nt}h(t)\dd t=0$ for $n=1,2,\dots$, then $g(u)=h(-\log u)$ satisfies $\int_0^1u^kg(u)\dd u=0$ for all $k\ge0$, so $g=0$ by Weierstrass' approximation theorem.
\end{proof}

\begin{theorem}[the generator of Brownian motion]\label{thm:bmgenerator}
Let $(P_t)$ be the Brownian semigroup in $\R^d$. Every $f\in C^2$ such that $f$, its first and second derivatives belong to $C_0$ is in $\mathcal D(A)$, and $Af=\frac12\Delta f$.
\end{theorem}

\begin{proof}
By It\^o's formula $f(x+B_t)=f(x)+\int_0^t\nabla f(x+B_s)\cdot\dd B_s+\frac12\int_0^t\Delta f(x+B_s)\dd s$, and the stochastic integral has a bounded integrand, hence mean $0$. So $P_tf-f=\int_0^tP_sg\dd s$ with
$g=\frac12\Delta f\in C_0$, and $\|\frac1t(P_tf-f)-g\|\le\frac1t\int_0^t\|P_sg-g\|\dd s\to0$.
\end{proof}

The class of functions in the theorem is not all of $\mathcal D(A)$, but it is enough to pin down the process. Stroock and Varadhan (1969) turned this into a method: a process \emph{solves the martingale
problem} for an operator $L$ if $f(X_t)-\int_0^tLf(X_s)\dd s$ is a martingale for every $f$ in a class of test functions, and one studies existence and uniqueness of solutions in law. For $\frac12\Delta$ uniqueness is
L\'evy's characterisation in disguise.

\begin{theorem}[the martingale problem for $\frac12\Delta$]\label{thm:martprob}
Let $X$ be a continuous adapted process in $\R^d$ with $X_0=0$, such that $f(X_t)-\frac12\int_0^t\Delta f(X_s)\dd s$ is a martingale for every $f\in C^2_c(\R^d)$. Then $X$ is a Brownian motion.
\end{theorem}

\begin{proof}
Let $\tau_k=\inf\{t:|X_t|\ge k\}$. For $f\in C^2(\R^d)$ choose $f_k\in C^2_c$ equal to $f$ on the ball $\bar B(0,k)$; then $f(X_{t\wedge\tau_k})-\frac12\int_0^{t\wedge\tau_k}\Delta f(X_s)\dd s$ is the stopped martingale for $f_k$. So
$f(X_t)-\frac12\int_0^t\Delta f(X_s)\dd s$ is a continuous local martingale for every $f\in C^2$. With $f(x)=x_i$: $X^i$ is a continuous local martingale. With $f(x)=x_ix_j$: $X^iX^j-\delta_{ij}t$ is a continuous local
martingale, so $\langle X^i,X^j\rangle_t=\delta_{ij}t$ (Proposition~\ref{prop:bracket}). L\'evy's characterisation (Theorem~\ref{thm:levychar}) concludes.
\end{proof}

\begin{remark}
The same proof shows that for $L=b\cdot\nabla+\frac12\sum_{i,j}a_{ij}\partial_{ij}$, with $a=\sigma\sigma^{\mathsf T}$, every continuous solution of the martingale problem has $X-\int b(X)\dd s$ a continuous local martingale with bracket
$\int a(X)\dd s$, and is then a solution of $\dd X=b(X)\dd t+\sigma(X)\dd W$ for some Brownian motion $W$, possibly on an enlarged probability space (a \emph{weak solution}). Uniqueness in law for the SDE is
equivalent to uniqueness for the martingale problem; when $b$ and $\sigma$ are Lipschitz it follows from Theorem~\ref{thm:sdeexist} and the Yamada--Watanabe principle that pathwise uniqueness implies uniqueness in law. Stroock and Varadhan proved
uniqueness when $a$ is merely continuous and uniformly positive definite (Stroock and Varadhan, 1979).
\end{remark}

\section{Poisson random measures}

\begin{definition}\label{def:prm}
Let $(E,\mathcal E,\mu)$ be a $\sigma$-finite measure space. A \emph{Poisson random measure} with intensity $\mu$ is a random measure $N$ on $E$ with values in $\{0,1,2,\dots\}\cup\{\infty\}$ such that, for pairwise disjoint
$A_1,\dots,A_k\in\mathcal E$, the variables $N(A_1),\dots,N(A_k)$ are independent, and $N(A)$ is Poisson$(\mu(A))$ (meaning $N(A)=\infty$ almost surely if $\mu(A)=\infty$).
\end{definition}

\begin{theorem}\label{thm:prmexist}
For every $\sigma$-finite $\mu$ there is a Poisson random measure with intensity $\mu$. If $\mu(E)<\infty$, one is $N=\sum_{k=1}^K\delta_{Y_k}$, where $K$ is Poisson$(\mu(E))$ and the $Y_k$ are independent with law
$\mu/\mu(E)$, independent of $K$.
\end{theorem}

\begin{proof}
If $\mu(E)<\infty$ and $A_1,\dots,A_k$ is a partition of $E$, then given $K=n$ the vector $(N(A_1),\dots,N(A_k))$ is multinomial with $n$ trials and probabilities $\mu(A_j)/\mu(E)$, so the colouring lemma gives the independence and the laws; disjoint sets that do
not cover $E$ are completed by their complement. In general write $E=\bigcup_jE_j$ with disjoint $E_j$ of finite measure, build independent Poisson random measures $N_j$ with intensities $\mu(\cdot\cap E_j)$, and
put $N=\sum_jN_j$. For disjoint $A_i$, the family $(N_j(A_i\cap E_j))_{i,j}$ is independent, and $N(A_i)=\sum_jN_j(A_i\cap E_j)$ is a sum of independent Poisson variables with means adding up to $\mu(A_i)$:
Poisson$(\mu(A_i))$, and almost surely infinite if $\mu(A_i)=\infty$ (because $P(N(A_i)\le m)$ is at most the probability that a Poisson variable of mean $\sum_{j\le J}\mu(A_i\cap E_j)$ is at most $m$, which tends to $0$ as $J\to\infty$).
\end{proof}

For a measurable $f$ we write $N(f)=\int f\dd N=\sum f(\text{atoms})$ when it makes sense.

\begin{proposition}[exponential formula]\label{prop:campbell}
Let $N$ be a Poisson random measure with intensity $\mu$.
\begin{enumerate}[label=(\alph*)]
\item For measurable $f\ge0$, $\mathbb Ee^{-N(f)}=\exp\bigl(-\int(1-e^{-f})\dd\mu\bigr)$ and $\mathbb EN(f)=\int f\dd\mu$.
\item If $\int(|f|\wedge1)\dd\mu<\infty$, then $N(|f|)<\infty$ almost surely and $\mathbb Ee^{iN(f)}=\exp\bigl(\int(e^{if}-1)\dd\mu\bigr)$.
\item If $f\in L^1(\mu)\cap L^2(\mu)$, then $\mathbb E\bigl(N(f)-\int f\dd\mu\bigr)^2=\int f^2\dd\mu$.
\end{enumerate}
\end{proposition}

\begin{proof}
(a) For $f=\sum_ja_j\mathbf 1_{A_j}$ simple with disjoint $A_j$ of finite measure, independence and $\mathbb Ee^{-aN(A)}=e^{-\mu(A)(1-e^{-a})}$ give the formula; a non-negative measurable $f$ is an increasing limit of such
functions (on a $\sigma$-finite space), and monotone convergence passes to the limit on both sides. Replacing $f$ by $sf$ and differentiating at $s=0$ (monotone convergence again) gives the mean.

(b) By (a), $\mathbb EN(|f|\mathbf 1_{\{|f|\le1\}})$ and $\mathbb EN(\{|f|>1\})$ are finite, so $N$ has finitely many atoms in $\{|f|>1\}$ and $N(|f|)<\infty$ almost surely. The characteristic function is computed for simple $f$ as in (a),
and the general case follows by dominated convergence, approximating $f$ by simple $f_n$ with $|f_n|\le|f|$, $f_n\to f$: $|e^{if_n}-1|\le2\wedge|f|$ is $\mu$-integrable.

(c) For simple $f$, $N(f)-\int f\dd\mu=\sum_ja_j(N(A_j)-\mu(A_j))$ is a sum of independent centred variables with variances $a_j^2\mu(A_j)$. For general $f$, take simple $f_n\to f$ in $L^1\cap L^2$ with
$|f_n|\le|f|$: by (a) applied to $|f-f_n|$, $N(f_n)\to N(f)$ in $L^1$, and the $N(f_n)-\int f_n\dd\mu$ form a Cauchy sequence in $L^2$ by the simple case, so the identity passes to the limit.
\end{proof}

\section{L\'evy processes}

\begin{definition}\label{def:levyproc}
A \emph{L\'evy process} is a real process $X$ with $X_0=0$, right-continuous paths with left limits (\emph{c\`adl\`ag}), and stationary independent increments: for $0\le t_0<\dots<t_k$, the increments
$X_{t_j}-X_{t_{j-1}}$ are independent, and the law of $X_{t+h}-X_t$ depends only on $h$.
\end{definition}

Brownian motion with drift, the Poisson process and the compound Poisson processes are L\'evy processes, and so is $(T_a)_{a\ge0}$, the hitting times of Brownian motion indexed by the level (Theorem~\ref{thm:halfstable}).

\begin{proposition}\label{prop:levyexp}
If $X$ is a L\'evy process, there is a unique triple $(b,\sigma^2,\Pi)$ such that $\mathbb Ee^{i\xi X_t}=e^{t\psi(\xi)}$ for all $t\ge0$ and $\xi\in\R$, with $\psi$ the L\'evy--Khintchine exponent of the triple. Moreover $P_tf(x)=\mathbb Ef(x+X_t)$ is a
Feller semigroup, and $X$ is a Markov process with this semigroup for its natural filtration.
\end{proposition}

\begin{proof}
Let $\varphi_t(\xi)=\mathbb Ee^{i\xi X_t}$. Independence and stationarity give $\varphi_{t+s}=\varphi_t\varphi_s$. Since $X_{t+h}-X_t$ has the law of $X_h$, which tends to $0$ in probability as $h\downarrow0$ by right-continuity,
$t\mapsto\varphi_t(\xi)$ is continuous. If $\varphi_1(\xi)=0$, then $\varphi_{1/n}(\xi)=0$ for all $n$, contradicting $\varphi_{1/n}(\xi)\to1$. Let $\psi$ be the continuous logarithm of $\varphi_1$. The function
$\varphi_{1/n}e^{-\psi/n}$ is continuous, equal to $1$ at $\xi=0$, and its $n$-th power is $1$: it takes values in the $n$-th roots of unity, so it is the constant $1$. Hence $\varphi_t=e^{t\psi}$ for rational $t$, and for all $t$ by
continuity. $X_1$ is the sum of the $n$ i.i.d.\ increments over $[\frac{k-1}n,\frac kn]$, so it is infinitely divisible, and L\'evy--Khintchine gives the triple. The Feller property is proved as for Brownian motion, using
$X_t\to0$ in probability as $t\to0$; the Markov property is the independence of $X_{t+s}-X_t$ from $(X_u)_{u\le t}$.
\end{proof}

\begin{proposition}[generator of a L\'evy process]\label{prop:levygen}
For $f$ in the Schwartz class $\mathcal S(\R)$, $f\in\mathcal D(A)$ and
\[
Af(x)=bf'(x)+\frac{\sigma^2}2f''(x)+\int\bigl(f(x+y)-f(x)-yf'(x)\mathbf 1_{\{|y|\le1\}}\bigr)\Pi(\mathrm dy).
\]
\end{proposition}

\begin{proof}
Write $f(x)=\frac1{2\pi}\int\hat f(\xi)e^{i\xi x}\dd\xi$ with $\hat f\in\mathcal S$. Then $P_tf(x)=\frac1{2\pi}\int\hat f(\xi)e^{t\psi(\xi)}e^{i\xi x}\dd\xi$ by Fubini, and
\[
\Bigl|\frac{P_tf(x)-f(x)}t-\frac1{2\pi}\int\hat f(\xi)\psi(\xi)e^{i\xi x}\dd\xi\Bigr|\le\frac1{2\pi}\int|\hat f(\xi)|\,\Bigl|\frac{e^{t\psi(\xi)}-1}t-\psi(\xi)\Bigr|\dd\xi .
\]
Since $\operatorname{Re}\psi\le0$, $|e^{z}-1-z|\le|z|^2/2$ and $|e^z-1|\le|z|$ for $\operatorname{Re}z\le0$; so the integrand is at most $2|\hat f||\psi|$ and tends to $0$ pointwise, and $|\psi(\xi)|\le C(1+\xi^2)$ (bound the integrand of
$\psi$ by $\frac{\xi^2x^2}2$ for $|x|\le1$ and by $2$ for $|x|>1$). Dominated convergence gives uniform convergence in $x$. Finally, $\frac1{2\pi}\int\hat f\psi e^{i\xi x}\dd\xi$ is the right side: the drift and Gaussian
terms are $bf'$ and $\frac{\sigma^2}2f''$, and for the integral term Fubini is allowed since $|e^{i\xi y}-1-i\xi y\mathbf 1_{\{|y|\le1\}}|\le(\xi^2y^2/2)\vee2$, which is integrable for $|\hat f(\xi)|\dd\xi\,\Pi(\mathrm dy)$.
\end{proof}

The generator reads like the L\'evy--Khintchine formula: drift, diffusion, and for each possible jump $y$, at rate $\Pi(\mathrm dy)$, the change $f(x+y)-f(x)$, corrected for the small jumps. The next theorem says that the process
itself is built in exactly this way.

\begin{theorem}[L\'evy--It\^o, 1934--1942]\label{thm:levyito}
\begin{enumerate}[label=(\alph*)]
\item Let $(b,\sigma^2,\Pi)$ be a triple, $W$ a Brownian motion and $N$ an independent Poisson random measure on $[0,\infty[\times\R^*$, where $\R^*=\R\setminus\{0\}$, with intensity $\mathrm dt\otimes\Pi$. For
$0<\varepsilon\le1$ let $D_\varepsilon=\{\varepsilon<|x|\le1\}$ and $D=\{|x|>1\}$, and put
\begin{gather*}
Y^\varepsilon_t=\int_{[0,t]\times D_\varepsilon}x\,N(\mathrm ds,\mathrm dx)-t\int_{D_\varepsilon}x\,\Pi(\mathrm dx),\\
Z_t=\int_{[0,t]\times D}x\,N(\mathrm ds,\mathrm dx).
\end{gather*}
Then, along a sequence $\varepsilon_k\downarrow0$, $Y^{\varepsilon_k}$ converges almost surely, uniformly on compact intervals, to a c\`adl\`ag process $Y$, and $X_t=bt+\sigma W_t+Y_t+Z_t$ is a L\'evy process with exponent $\psi$.
Its jumps are the atoms of $N$: $N=\sum_{s:\,\Delta X_s\ne0}\delta_{(s,\Delta X_s)}$.
\item Conversely, let $X$ be a L\'evy process with triple $(b,\sigma^2,\Pi)$, and $N^X=\sum_{s:\,\Delta X_s\ne0}\delta_{(s,\Delta X_s)}$ its jump measure. Then $N^X$ is a Poisson random measure with intensity $\mathrm dt\otimes\Pi$, the process
$X_t-bt-Y_t-Z_t$, with $Y,Z$ built from $N^X$ as in (a), is $\sigma W$ for a Brownian motion $W$ independent of $N^X$, and the convergence of (a) holds.
\end{enumerate}
\end{theorem}

\begin{proof}
(a) \emph{Large jumps.} $\Pi(|x|>1)<\infty$, so $N$ has finitely many atoms in $[0,t]\times\{|x|>1\}$ (Proposition~\ref{prop:campbell}(a)), and $Z$ is a piecewise constant c\`adl\`ag process, a compound Poisson process. The
same holds for $Y^\varepsilon$, plus a linear drift. By Proposition~\ref{prop:campbell}(b), for $s<t$,
\begin{gather*}
\mathbb Ee^{i\xi(Z_t-Z_s)}=\exp\Bigl((t-s)\int_D(e^{i\xi x}-1)\Pi(\mathrm dx)\Bigr),\\
\mathbb Ee^{i\xi(Y^\varepsilon_t-Y^\varepsilon_s)}=e^{(t-s)\psi_\varepsilon(\xi)},\\\psi_\varepsilon(\xi)=\int_{D_\varepsilon}(e^{i\xi x}-1-i\xi x)\Pi(\mathrm dx),
\end{gather*}
and increments over disjoint time intervals, or of $Z$ and $Y^\varepsilon$, are built from $N$ on disjoint sets, hence independent.

\emph{Small jumps.} $Y^\varepsilon$ has independent centred increments, so it is a martingale for the filtration $\mathcal G_t$ generated by $W$ and the restriction of $N$ to $[0,t]\times\R$, and by Proposition~\ref{prop:campbell}(c),
for $\varepsilon'<\varepsilon$, $\mathbb E(Y^{\varepsilon'}_t-Y^\varepsilon_t)^2=t\int_{\varepsilon'<|x|\le\varepsilon}x^2\Pi(\mathrm dx)$. Doob's inequality, which holds for c\`adl\`ag martingales with the same proof as Theorem~\ref{thm:doobcont}
(apply the discrete inequality on finite sets increasing to a countable dense set), gives
\[
\mathbb E\sup_{s\le t}|Y^{\varepsilon'}_s-Y^\varepsilon_s|^2\le4t\int_{\varepsilon'<|x|\le\varepsilon}x^2\Pi(\mathrm dx).
\]
The right side tends to $0$ as $\varepsilon\to0$, since $\int_{|x|\le1}x^2\Pi<\infty$. Choose $\varepsilon_k\downarrow0$ with $\int_{|x|\le\varepsilon_k}x^2\Pi\le4^{-k}$: then $\sum_k\mathbb E\sup_{s\le t}|Y^{\varepsilon_{k+1}}_s-Y^{\varepsilon_k}_s|\le\sum_k2\sqrt t\,2^{-k}<\infty$,
so $\sum_k\sup_{s\le t}|Y^{\varepsilon_{k+1}}_s-Y^{\varepsilon_k}_s|<\infty$ almost surely, and $Y^{\varepsilon_k}$ converges uniformly on $[0,t]$ to a c\`adl\`ag limit $Y$ (a uniform limit of c\`adl\`ag functions is c\`adl\`ag). Independence and stationarity of increments pass to the limit, and
$\mathbb Ee^{i\xi Y_t}=\exp(t\int_{|x|\le1}(e^{i\xi x}-1-i\xi x)\Pi(\mathrm dx))$ by dominated convergence ($|e^{i\xi x}-1-i\xi x|\le\xi^2x^2/2$). Multiplying the characteristic functions of the independent pieces gives $e^{t\psi}$.

\emph{Jumps.} Almost surely no two atoms of $N$ have the same time coordinate: on $[0,T]\times\{|x|>\varepsilon\}$ the intensity is finite, and by Theorem~\ref{thm:prmexist} the atoms there are finitely many independent points
whose time coordinates are uniform on $[0,T]$, hence distinct. So each $Y^\varepsilon+Z$ jumps exactly at the atoms $(s,x)$ with $|x|>\varepsilon$, by $x$. At a time $s$, the uniform limit $Y+Z$ has jump $\lim_k\Delta(Y^{\varepsilon_k}+Z)_s$, which is $x$ if $(s,x)$ is an atom and
$0$ otherwise. $W$ is continuous.

(b) Let $X'$ be the process constructed in (a) from the same triple. $X$ and $X'$ have the same finite-dimensional laws, since both have independent stationary increments with $\mathbb Ee^{i\xi X_t}=e^{t\psi(\xi)}$ (Proposition~\ref{prop:levyexp}).
A c\`adl\`ag path is determined by its values at rational times, and so are its jumps: for $a<b$ with $0\notin[a,b]$ and $t$ rational, the number of jumps in $[0,t]$ with size in $]a,b[$ is a limit of counts of rational
increments, and the processes $Y^\varepsilon,Z$ built from the jump measure, the supremum over $[0,t]$ of $|X-bt-Y^{\varepsilon}-Z-(X'-\dots)|$, and so on, are measurable functions of the values at rational times. Therefore the
joint law of $(X,N^X,(Y^{\varepsilon_k})_k,Z)$ is the joint law of the same objects for $X'$. For $X'$, part (a) shows that $N^{X'}=N$ is a Poisson random measure with intensity $\mathrm dt\otimes\Pi$, that $Y^{\varepsilon_k}$
converges uniformly on compacts, and that $X'-bt-Y-Z=\sigma W$ with $W$ a Brownian motion independent of $N$. These are properties of the joint law, so they hold for $X$.
\end{proof}

\begin{figure}[ht]
\centering
\fig{ch07-subordinator}
\caption{Top: a Brownian path $B$ and its running maximum $M$ on $[0,1]$. Bottom: the hitting times $a\mapsto T_a$, $0\le a\le M_1$, a $\frac12$-stable subordinator. Each jump of $T$ is the length of an interval on which $M$
is constant: a flat stretch of the top picture becomes a vertical step of the bottom one.}
\label{fig:subordinator}
\end{figure}

\section{Subordinators and the clock of Brownian excursions}

A \emph{subordinator} is a L\'evy process with non-decreasing paths: a random clock.

\begin{corollary}\label{cor:subordinator}
A L\'evy process $X$ is a subordinator if and only if $\sigma^2=0$, $\Pi$ is carried by $]0,\infty[$, $\int(x\wedge1)\Pi(\mathrm dx)<\infty$ and $d=b-\int_{]0,1]}x\,\Pi(\mathrm dx)\ge0$. Then
$X_t=dt+\sum_{s\le t}\Delta X_s$ and $\mathbb Ee^{-\lambda X_t}=e^{-t\Phi(\lambda)}$ with $\Phi(\lambda)=d\lambda+\int(1-e^{-\lambda x})\Pi(\mathrm dx)$, the \emph{Laplace exponent}.
\end{corollary}

\begin{proof}
If $X$ is non-decreasing, its jumps are positive, so $\Pi(]-\infty,0[)=0$ by Theorem~\ref{thm:levyito}(b), and $J_t=\sum_{s\le t}\Delta X_s\le X_t<\infty$. By Proposition~\ref{prop:campbell}(a),
$\mathbb Ee^{-J_t}=\exp(-t\int(1-e^{-x})\Pi(\mathrm dx))$ is positive, so $\int(1-e^{-x})\Pi<\infty$, which is $\int(x\wedge1)\Pi<\infty$. Then $Y_t=J^{\le1}_t-t\int_{]0,1]}x\Pi$, where $J^{\le1}$ is the sum of the jumps of size at most $1$
(the compensated sums converge to this, by monotone convergence and Proposition~\ref{prop:campbell}(a)), and Theorem~\ref{thm:levyito}(b) gives $X_t=J_t+dt+\sigma W_t$. The process $X_t-J_t=dt+\sigma W_t$ is non-decreasing (a jump of
$X$ is exactly matched in $J$), which forces $\sigma=0$ (Brownian motion has infinite variation) and $d\ge0$. Conversely such an $X$ is non-decreasing. The Laplace transform follows from Proposition~\ref{prop:campbell}(a).
\end{proof}

\begin{theorem}[the $\frac12$-stable subordinator]\label{thm:halfstable}
Let $B$ be a Brownian motion and $T_a=\inf\{t:B_t>a\}$ for $a\ge0$. Then $(T_a)_{a\ge0}$ is a subordinator with Laplace exponent $\Phi(\lambda)=\sqrt{2\lambda}$, no drift, and L\'evy measure
\[
\Pi(\mathrm dx)=\frac{\mathrm dx}{\sqrt{2\pi}\,x^{3/2}}\qquad(x>0).
\]
It is $\frac12$-stable: $(T_{ca})_{a\ge0}$ has the law of $(c^2T_a)_{a\ge0}$.
\end{theorem}

\begin{proof}
$T_0=0$ almost surely (Blumenthal, Theorem~\ref{thm:blumenthal}). $a\mapsto T_a$ is non-decreasing and right-continuous: $\{T_a<t\}=\bigcup_n\{T_{a+1/n}<t\}$ follows from $\{T_a<t\}=\{\max_{s\le t'}B_s>a\text{ for some }t'<t\}$.
For $a<a'$, $T_{a'}-T_a$ is the first time $B_{T_a+\cdot}-B_{T_a}$ exceeds $a'-a$ (since $B_{T_a}=a$ by continuity); by the strong Markov property (Proposition~\ref{prop:strongmarkovd}) it is independent of $\mathcal F_{T_a}$, which contains
$(T_u)_{u\le a}$, and has the law of $T_{a'-a}$. So $(T_a)$ is a subordinator, with $\mathbb Ee^{-\lambda T_a}=e^{-a\sqrt{2\lambda}}$ (Volume~10, Theorem~5.10). Integrating by parts,
$\int_0^\infty(1-e^{-\lambda x})x^{-3/2}\dd x=\int_0^\infty2x^{-1/2}\lambda e^{-\lambda x}\dd x=2\sqrt{\pi\lambda}$, so the stated $\Pi$ and $d=0$ give $\Phi(\lambda)=\sqrt{2\lambda}$; the Laplace transform determines the law of $T_1$, which determines its triple (Volume~10, Theorem~5.3(c)), hence $d$ and $\Pi$. Scaling: $(c^{-1}B_{c^2t})$ is a Brownian motion whose hitting time of $a$ is $c^{-2}T_{ca}$.
\end{proof}

The jumps of $T$ are the lengths of the intervals on which the running maximum $M$ is constant: the excursions of $B$ below its maximum (Figure~\ref{fig:subordinator}). By Theorem~\ref{thm:levyito}(b), the number of
these excursions longer than $x$ completed before $B$ first exceeds $a$ is Poisson with mean $a\,\Pi(]x,\infty[)=a\sqrt{\frac2{\pi x}}$. Infinitely many short ones, a few long ones, independent of each other: this is the
structure that It\^o turned in 1970 into his theory of \emph{excursions}, a Poisson random measure whose atoms are whole excursions.

\begin{corollary}[the inverse local time]\label{cor:invlocal}
Let $L$ be the local time at $0$ of a Brownian motion and $\tau_\ell=\inf\{t:L_t>\ell\}$. Then $(\tau_\ell)_{\ell\ge0}$ is a $\frac12$-stable subordinator with Laplace exponent $\sqrt{2\lambda}$.
\end{corollary}

\begin{proof}
By L\'evy's theorem (Theorem~\ref{thm:levyMB}), the process $L$ has the law of the running maximum $M$, and $\tau_\ell$ is the same functional of $L$ as $T_\ell=\inf\{t:M_t>\ell\}$ is of $M$.
\end{proof}

\begin{beauty}
Local time is a clock that only ticks when the Brownian particle is at $0$; turned inside out, $\tau_\ell$ is a clock indexed by local time, and it runs by jumps, one for each excursion away from $0$. The law of
these jumps, $\frac{\mathrm dx}{\sqrt{2\pi}x^{3/2}}$, is forced by scaling alone: a Brownian excursion of length $x$ and one of length $c^2x$ differ by the factor $c$ in height, and only the exponent $\frac32$ is compatible.
One can show that the zero set of Brownian motion, of dimension $\frac12$ (Corollary~\ref{cor:zerodim}), is the closure of the range of this subordinator.
\end{beauty}

\begin{think}
The $\frac12$-stable subordinator has infinitely many jumps in every interval, yet its value $T_a$ is finite. Which property of $\Pi$ makes both true at once? What would a L\'evy measure $\frac{\mathrm dx}{x^2}$ on $]0,\infty[$
give?
\end{think}

\begin{remark}
Stable processes of every index $\alpha\in]0,2[$ arise the same way: $\Pi(\mathrm dx)=c_\pm|x|^{-1-\alpha}\mathrm dx$ on each half-line gives, with a suitable drift, a process with $X_{ct}\overset{d}=c^{1/\alpha}X_t$, and Brownian motion evaluated at
an independent $\frac12$-stable clock is the Cauchy process, of index $1$ (Exercise~\ref{ex:cauchy}). Bertoin's book is the reference.
\end{remark}

\section{Exercises}

\begin{exercise}[$\star$]\label{ex:poissongen}
Let $N$ be a Poisson process of rate $\lambda$. Show that it is a L\'evy process with triple $(\lambda,0,\lambda\delta_1)$, and that its generator is $Af(x)=\lambda(f(x+1)-f(x))$ on all of $C_0$. Write Dynkin's martingale for
$f$ and check it directly for $f(x)=x$ (locally).
\end{exercise}

\begin{exercise}[$\star$]\label{ex:driftbm}
Show that $X_t=B_t+\mu t$ is a Feller process with generator $\frac12f''+\mu f'$ on $C^2$ functions with $f,f',f''\in C_0$. Using Dynkin's formula with $f(x)=e^{-2\mu x}$ (suitably modified outside an interval), recover the probability that $X$
reaches $-a$ before $c$ (Chapter~3).
\end{exercise}

\begin{exercise}[$\star$]\label{ex:gammasub}
Let $\Pi(\mathrm dx)=x^{-1}e^{-x}\mathrm dx$ on $]0,\infty[$. Show that the subordinator with this L\'evy measure and no drift has $X_t$ with the Gamma$(t,1)$ law. It has infinitely many jumps on every interval; what is the
expected number of jumps larger than $x$ on $[0,1]$?
\end{exercise}

\begin{exercise}[$\star\star$]\label{ex:mehler}
Let $P_tf(x)=\mathbb Ef(e^{-t}x+\sqrt{1-e^{-2t}}Z)$, $Z\sim N(0,1)$ (the Ornstein--Uhlenbeck semigroup, Proposition~\ref{prop:ou} with $\theta=1$, $\sigma=\sqrt2$). Show that $(P_t)$ is a Feller semigroup and that
$Af=f''-xf'$ for $f\in C^2_c$. Show that $\int P_tf\,\mathrm d\gamma=\int f\dd\gamma$ for the standard Gaussian measure $\gamma$.
\end{exercise}

\begin{exercise}[$\star\star$]\label{ex:stable}
Let $0<\alpha<2$ and $\Pi(\mathrm dx)=|x|^{-1-\alpha}\mathrm dx$ on $\R\setminus\{0\}$, $b=\sigma=0$. Show that $\psi(\xi)=-k_\alpha|\xi|^\alpha$ with $k_\alpha=\int(1-\cos u)|u|^{-1-\alpha}\dd u$, and that $(X_{ct})_t$ has the law of $(c^{1/\alpha}X_t)_t$. For which
$\alpha$ is $\mathbb E|X_1|<\infty$?
\end{exercise}

\begin{exercise}[$\star\star$]\label{ex:cauchy}
Let $T$ be the $\frac12$-stable subordinator of Theorem~\ref{thm:halfstable} and $W$ an independent Brownian motion. Show that $C_a=W_{T_a}$ is a L\'evy process with $\mathbb Ee^{i\xi C_a}=e^{-a|\xi|}$: a Cauchy process. Relate this to
Exercise~\ref{ex:halfplane}.
\end{exercise}

\begin{exercise}[$\star\star$]\label{ex:compound}
Let $X$ be a compound Poisson process: $X_t=\sum_{k\le N_t}\xi_k$, with $N$ a Poisson process of rate $\lambda$ and $\xi_k$ i.i.d.\ with law $\nu$, independent of $N$. Find its triple, and show that if $\nu$ has a finite mean $m$, then $X_t-\lambda mt$
is a martingale. With $\nu$ on $]0,\infty[$, $X$ is the total claim of an insurance company; show that $u+ct-X_t$ (capital $u$, premium rate $c>\lambda m$) satisfies $\mathbb E[\text{capital at }t]\to\infty$, and discuss why ruin can still happen.
\end{exercise}

\begin{exercise}[$\star\star\star$]\label{ex:excursions}
(a) Show that the number of jumps of $(T_a)_{a\le1}$ larger than $x$ is Poisson with mean $\sqrt{2/(\pi x)}$, and deduce the law of the largest jump $J$: $P(J\le x)=e^{-\sqrt{2/(\pi x)}}$.
(b) Interpret $J$ as the longest excursion of $B$ below its running maximum before $B$ first exceeds $1$, and show that $\mathbb E\sqrt{J}=\infty$ but $\mathbb EJ^{\beta}<\infty$ for $\beta<\frac12$.
(c) Show that $T_1=\sum(\text{jumps})$ and that $P(T_1>t)$ and $P(J>t)$ are both of order $t^{-1/2}$ as $t\to\infty$: a large $T_1$ is made of one large jump.
\end{exercise}

\section{Hints}

\begin{hint}{ex:poissongen}
$\mathbb Ee^{i\xi N_t}=e^{\lambda t(e^{i\xi}-1)}$. For the generator, $P_tf(x)=\sum_ke^{-\lambda t}\frac{(\lambda t)^k}{k!}f(x+k)$.
\end{hint}

\begin{hint}{ex:driftbm}
$e^{-2\mu x}$ is annihilated by $\frac12f''+\mu f'$. Modify it outside $[-a-1,c+1]$ to lie in $C_0$, and stop at the exit time of $]-a,c[$.
\end{hint}

\begin{hint}{ex:gammasub}
$\int_0^\infty(1-e^{-\lambda x})x^{-1}e^{-x}\dd x=\log(1+\lambda)$ (differentiate in $\lambda$).
\end{hint}

\begin{hint}{ex:mehler}
Write $P_tf(x)$ as an integral against a Gaussian density. For the generator, It\^o's formula for the OU process, or a Taylor expansion. For invariance, if $X_0\sim\gamma$ then $X_t\sim\gamma$.
\end{hint}

\begin{hint}{ex:stable}
The integrand of $\psi$ is odd in its imaginary part. Substitute $u=|\xi|x$. Scaling: compare the exponents of $X_{ct}$ and $c^{1/\alpha}X_t$.
\end{hint}

\begin{hint}{ex:cauchy}
Condition on $T$: $\mathbb E[e^{i\xi W_{T_a}}\mid T]=e^{-\xi^2T_a/2}$.
\end{hint}

\begin{hint}{ex:compound}
Proposition~\ref{prop:campbell}(b) with $\mu=\lambda\,\mathrm dt\otimes\nu$. For ruin, the capital can drop below $0$ at a jump.
\end{hint}

\begin{hint}{ex:excursions}
(a) Theorem~\ref{thm:levyito}(b) for the set $[0,1]\times]x,\infty[$. (b) $\mathbb EJ^\beta=\int_0^\infty\beta x^{\beta-1}P(J>x)\dd x$ and $1-e^{-y}\sim y$. (c) $P(T_1>t)=P(M_t<1)$, and the reflection principle.
\end{hint}

\section{Solutions}

\begin{solution}{ex:poissongen}
\idea A Poisson process is the purest jump process: drift $0$ except for the bookkeeping of the truncation, no Gaussian part, and jumps of size $1$ at rate $\lambda$.
\writeup Independent Poisson increments give stationarity and independence, and $\log\mathbb Ee^{i\xi N_1}=\lambda(e^{i\xi}-1)=i\lambda\xi+\int(e^{i\xi x}-1-i\xi x\mathbf 1_{\{|x|\le1\}})\lambda\delta_1(\mathrm dx)$, so the triple is
$(\lambda,0,\lambda\delta_1)$. For $f\in C_0$, $\frac1t(P_tf-f)(x)=\frac{e^{-\lambda t}-1}tf(x)+\lambda e^{-\lambda t}f(x+1)+O(t)\|f\|$, uniformly in $x$, so $Af=\lambda(f(\cdot+1)-f)$ for every $f\in C_0$. Dynkin's martingale is
$f(N_t)-f(0)-\lambda\int_0^t(f(N_s+1)-f(N_s))\dd s$; for $f(x)=x$ (which is not in $C_0$, but the computation is the same after stopping when $N$ reaches a level), it is $N_t-\lambda t$, the compensated Poisson process.
\end{solution}

\begin{solution}{ex:driftbm}
\idea Adding a drift adds a first-order term to the generator, and the function killed by the generator turns the exit problem into a linear equation.
\writeup The Feller property and the generator follow as in Theorem~\ref{thm:bmgenerator}, with It\^o's formula for $f(x+B_t+\mu t)$, which has the extra term $\int\mu f'\dd s$. Let $h(x)=e^{-2\mu x}$, so $\frac12h''+\mu h'=0$, and $g\in C^2$
with $g,g',g''\in C_0$ equal to $h$ on $[-a,c]$. With $\tau$ the exit time of $]-a,c[$ (finite with finite mean, since $X$ exits by Brownian fluctuations alone) and $t\wedge\tau$, Dynkin's formula gives $\mathbb Eh(X_{t\wedge\tau})=1$, and
letting $t\to\infty$, $p\,e^{2\mu a}+(1-p)e^{-2\mu c}=1$ with $p=P(X_\tau=-a)$. So $p=\frac{1-e^{-2\mu c}}{e^{2\mu a}-e^{-2\mu c}}$, as in Exercise~\ref{ex:driftruin}.
\end{solution}

\begin{solution}{ex:gammasub}
\idea A Gamma law is infinitely divisible with Laplace transform $(1+\lambda)^{-t}$; reading $\log(1+\lambda)$ as $\int(1-e^{-\lambda x})\Pi(\mathrm dx)$ produces the L\'evy measure.
\writeup $\int(x\wedge1)x^{-1}e^{-x}\dd x<\infty$, so Corollary~\ref{cor:subordinator} applies, with $\Phi(\lambda)=\int_0^\infty(1-e^{-\lambda x})x^{-1}e^{-x}\dd x$. Then $\Phi(0)=0$ and $\Phi'(\lambda)=\int_0^\infty e^{-(1+\lambda)x}\dd x=\frac1{1+\lambda}$, so
$\Phi(\lambda)=\log(1+\lambda)$ and $\mathbb Ee^{-\lambda X_t}=(1+\lambda)^{-t}$, the Laplace transform of Gamma$(t,1)$. The number of jumps larger than $x$ on $[0,1]$ is Poisson with mean $\int_x^\infty y^{-1}e^{-y}\dd y=E_1(x)\sim\log\frac1x$ as
$x\to0$: infinitely many small jumps, but logarithmically few.
\end{solution}

\begin{solution}{ex:mehler}
\idea The Ornstein--Uhlenbeck process pulls towards $0$ with force $-x$ and diffuses with $f''$; the Gaussian measure is its equilibrium.
\writeup $P_tf(x)=\int f(e^{-t}x+\sqrt{1-e^{-2t}}z)\gamma(\mathrm dz)$: by dominated convergence $P_tf\in C_0$ when $f\in C_0$, and $\|P_tf-f\|\le\sup_x\mathbb E|f(e^{-t}x+\sqrt{1-e^{-2t}}Z)-f(x)|\to0$ by uniform continuity of $f$, together with $\sup_x|f(e^{-t}x)-f(x)|\to0$ for $f\in C_0$ (split $|x|\le R$ and $|x|>R$). Positivity is clear and $P_tP_s=P_{t+s}$ is the Markov property of Proposition~\ref{prop:ou}
(Theorem~\ref{thm:sdemarkov}). For $f\in C^2_c$, Exercise~\ref{ex:sdegenerator} with $b(x)=-x$, $\sigma=\sqrt2$ gives pointwise convergence to $f''-xf'$; uniformity holds because $f$ has compact support and
$\frac1t\mathbb E|X^x_s-x|^2$ is bounded on compacts while for large $|x|$ both sides are small. Invariance: if $X_0\sim\gamma$ is independent of $Z$, then $e^{-t}X_0+\sqrt{1-e^{-2t}}Z\sim\gamma$, so $\int P_tf\dd\gamma=\int f\dd\gamma$.
\end{solution}

\begin{solution}{ex:stable}
\idea The L\'evy measure $|x|^{-1-\alpha}\mathrm dx$ is invariant, up to the factor $c^{\alpha}$, under $x\mapsto cx$: so the process is invariant under the matching change of space and time scales.
\writeup $\int(x^2\wedge1)|x|^{-1-\alpha}\dd x<\infty$ exactly when $0<\alpha<2$. By symmetry of $\Pi$ the terms $i\sin$ and $-i\xi x\mathbf 1$ integrate to $0$, so $\psi(\xi)=-\int(1-\cos\xi x)|x|^{-1-\alpha}\dd x=-|\xi|^\alpha k_\alpha$
after $u=|\xi|x$. Then $\mathbb Ee^{i\xi X_{ct}}=e^{-ctk_\alpha|\xi|^\alpha}=\mathbb Ee^{i(c^{1/\alpha}\xi)X_t}$, and since both processes have independent stationary increments, their finite-dimensional laws agree. $\mathbb E|X_1|<\infty$
iff $\int_{|x|>1}|x|\Pi(\mathrm dx)<\infty$ (the large jumps form an independent compound Poisson process, and the rest has all moments), that is iff $\alpha>1$.
\end{solution}

\begin{solution}{ex:cauchy}
\idea Running a Brownian motion on an independent random clock is subordination; the Gaussian characteristic function turns into the Laplace transform of the clock.
\writeup Given $T$, the increments $W_{T_{a_j}}-W_{T_{a_{j-1}}}$ are independent centred Gaussians with variances $T_{a_j}-T_{a_{j-1}}$, which are themselves independent and stationary; so $C$ has independent stationary increments, and
$\mathbb Ee^{i\xi C_a}=\mathbb Ee^{-\xi^2T_a/2}=e^{-a\sqrt{\xi^2}}=e^{-a|\xi|}$. $C$ is c\`adl\`ag since $T$ is and $W$ is continuous. This is the law of Exercise~\ref{ex:halfplane}: a planar Brownian motion $(W,\beta)$ started at height
$a$ hits the axis at time $T_a$ of $\beta$ (reflected), and its abscissa there is $W_{T_a}$, Cauchy with scale $a$; as $a$ grows, these hitting points form the Cauchy process.
\end{solution}

\begin{solution}{ex:compound}
\idea A compound Poisson process is a Poisson random measure on time $\times$ marks, summed; martingales come from subtracting the mean rate.
\writeup $N^X=\sum_k\delta_{(T_k,\xi_k)}$ is a Poisson random measure with intensity $\lambda\,\mathrm dt\otimes\nu$ (marking, as in Theorem~\ref{thm:prmexist}), so $\mathbb Ee^{i\xi X_t}=\exp(t\lambda\int(e^{i\xi x}-1)\nu(\mathrm dx))$ and the triple is
$(\lambda\int_{|x|\le1}x\,\nu(\mathrm dx),0,\lambda\nu)$. With a finite mean, $\mathbb EX_t=\lambda mt$ and independent increments make $X_t-\lambda mt$ a martingale. The expected capital $u+(c-\lambda m)t$ tends to $\infty$, but a large claim
early on (or an unlucky run) can take the capital below $0$; this is the Cram\'er--Lundberg model, in which the ruin probability is positive and decays exponentially in $u$ when $\nu$ has exponential moments (Asmussen, \emph{Ruin
Probabilities}, 2000).
\end{solution}

\begin{solution}{ex:excursions}
\idea A $\frac12$-stable subordinator is the sum of its jumps, and its tail comes from the single largest one, as for sums of heavy-tailed variables.
\writeup (a) By Theorem~\ref{thm:levyito}(b) and Theorem~\ref{thm:halfstable}, $N^T([0,1]\times]x,\infty[)$ is Poisson with mean $\int_x^\infty\frac{\mathrm dy}{\sqrt{2\pi}y^{3/2}}=\sqrt{\frac2{\pi x}}$, and $\{J\le x\}$ is the event that it vanishes.
(b) The jumps of $T$ on $[0,1]$ are the lengths of the maximal intervals before $T_1$ on which $M$ is constant. $P(J>x)=1-e^{-\sqrt{2/(\pi x)}}\sim\sqrt{2/(\pi x)}$ as $x\to\infty$, so $\mathbb EJ^\beta=\int_0^\infty\beta x^{\beta-1}P(J>x)\dd x$ is finite iff
$\beta-1-\frac12<-1$, that is $\beta<\frac12$. (c) $T$ has no drift, so $T_1$ is the sum of its jumps (Corollary~\ref{cor:subordinator}). $P(T_1>t)=P(M_t<1)=P(|B_t|<1)\sim\sqrt{\frac2{\pi t}}$ by the reflection principle, the same order
as $P(J>t)$, and in fact the same constant: the event $T_1>t$ is, to first order, the event that one jump exceeds $t$.
\end{solution}

\section*{Sources and further reading}

\begin{itemize}
\item \emph{History.} P.~L\'evy, ``Sur les int\'egrales dont les \'el\'ements sont des variables al\'eatoires ind\'ependantes'', \emph{Annali della Scuola Normale Superiore di Pisa} (2) 3 (1934), 337--366. K.~It\^o, ``On stochastic
processes (I) (Infinitely divisible laws of probability)'', \emph{Japanese Journal of Mathematics} 18 (1942), 261--301, and ``Poisson point processes attached to Markov processes'', \emph{Proceedings of the Sixth Berkeley
Symposium} 3 (1970), 225--239 (excursions). E.~Hille, \emph{Functional Analysis and Semi-Groups}, American Mathematical Society, 1948; K.~Yosida,
``On the differentiability and the representation of one-parameter semi-group of linear operators'', \emph{Journal of the Mathematical Society of Japan} 1 (1948), 15--21. W.~Feller, ``The parabolic differential equations
and the associated semi-groups of transformations'', \emph{Annals of Mathematics} 55 (1952), 468--519. D.~W.~Stroock and S.~R.~S.~Varadhan, ``Diffusion processes with continuous coefficients'', \emph{Communications on Pure
and Applied Mathematics} 22 (1969), 345--400 and 479--530, and \emph{Multidimensional Diffusion Processes}, Springer, 1979.
\item J.~Bertoin, \emph{L\'evy Processes}, Cambridge University Press, 1996. Subordinators, stable processes, fluctuation theory and excursions.
\item K.~Sato, \emph{L\'evy Processes and Infinitely Divisible Distributions}, Cambridge University Press, 1999. The L\'evy--It\^o decomposition with full details.
\item A.~E.~Kyprianou, \emph{Fluctuations of L\'evy Processes with Applications}, 2nd edition, Springer, 2014. A readable path to the applications, including insurance and finance.
\item S.~N.~Ethier and T.~G.~Kurtz, \emph{Markov Processes: Characterization and Convergence}, Wiley, 1986. Semigroups, generators and martingale problems.
\end{itemize}
